% ------------------------------------------------------------------
%  Capitolo 26 — Formare chi insegnerà
%  Parte VI
%  Ricerca di riferimento: ricerca 02 §8; 04; 01 §9.7, §14; ricerca 09 (da fare)
%  Voci bibliografiche principali: dekker2012, bei2023, oecd2025talis, pashler2008, kirschner2006
%  Epigrafe: ✔ verificata sul testo (vedi ricerca 03 e registro, ricerca 00)
%  Scheda del capitolo: ricerca/06_indice-ragionato.md, capitolo 26
% ------------------------------------------------------------------
\chapter{Formare chi insegnerà}
\label{cap:formatori}

\epigrafe[Grande didattica, che presenta l'arte universale di insegnare tutto a tutti]{Didactica magna, universale omnes omnia docendi artificium exhibens}{Comenio, Didactica magna (1657), frontespizio}

\iniziale{Q}{uesto capitolo} \segnaposto{apertura: da scrivere — vedi la scheda nell'indice ragionato}.

\section{Chi forma gli insegnanti}

\segnaposto{da scrivere}

\section{Che cosa imparano i futuri insegnanti, e che cosa manca}

\segnaposto{da scrivere}

\section{I neuromiti nella formazione}

\segnaposto{da scrivere}

\section{Un curricolo minimo di scienza dell'apprendimento}

\segnaposto{da scrivere}

\section{Dall'università alla scuola: la seconda fase}

\segnaposto{da scrivere}

\begin{daricordare}
\segnaposto{il principio del capitolo in due righe}
\end{daricordare}

\begin{domande}
  \item \segnaposto{domanda di recupero su questo capitolo}
  \item \segnaposto{domanda di applicazione a una materia che studi}
  \item \segnaposto{domanda di ripresa su un capitolo precedente}
\end{domande}

\begin{sintesi}
  \item \segnaposto{punto di sintesi}
\end{sintesi}

\finecapitolo
